\chapter{Transcutaneous Electrical Nerve Stimulation (TENS)}
\index{Transcutaneous Electrical Nerve Stimulation (TENS)}
\medskip
\noindent\textit{Educational reference; always seek professional advice for individual care.}

\section*{Definition and Overview}
Transcutaneous electrical nerve stimulation (TENS) uses a small battery-powered device to send mild electrical pulses through electrodes placed on the skin. It may reduce pain for some people with musculoskeletal, neuropathic, postoperative, or labour pain. It is an adjunct to diagnosis and treatment, not a cure for the underlying cause. Evidence varies by condition and device settings.

\section*{How It May Help}
Electrical stimulation may reduce pain signalling in the nervous system and may activate the body's own pain-modulating pathways. Settings differ in frequency, pulse width, intensity, and treatment time. A physiotherapist or other trained clinician can help with electrode placement and safe use.

\section*{Using TENS}
Place electrodes on intact skin near, but not directly over, the painful area as instructed. Start with a comfortable intensity and follow the device instructions. Do not use it while driving, sleeping, bathing, swimming, or operating machinery. Stop if it causes skin damage, worsening pain, dizziness, palpitations, or other concerning symptoms.

\section*{When to Seek Medical Care}
New severe pain, weakness, numbness, fever, unexplained weight loss, chest pain, or loss of bladder or bowel control needs medical assessment rather than self-treatment. \textbf{Contact local emergency services} for severe breathlessness, collapse, sudden neurological symptoms, or suspected serious injury.

\section*{Safety and Contraindications}
Do not place electrodes over a pacemaker or implanted cardiac defibrillator, the front or side of the neck, the head, eyes, or broken or infected skin. Avoid placing them across the chest unless specifically advised by a specialist. Ask a clinician before use in pregnancy, epilepsy, reduced sensation, bleeding problems, cancer, or significant heart disease. Do not use TENS to replace prescribed treatment.

\section*{Outlook}
Some people obtain short-term relief; others notice little benefit. Keep a brief record of pain, activity, and side effects. Review continued use with a clinician if there is no meaningful improvement after a reasonable trial.

\section*{Sources and Further Reading}
\begin{itemize}
    \item NHS. \textit{Transcutaneous electrical nerve stimulation (TENS).} https://www.nhs.uk/conditions/tens/
    \item Cochrane. \textit{TENS for pain management.} https://www.cochrane.org/
    \item Cleveland Clinic. \textit{TENS unit.} https://my.clevelandclinic.org/
\end{itemize}
